\documentclass[11pt]{article}
\usepackage[margin=25mm]{geometry}
\usepackage{fontspec}\setmainfont{Latin Modern Roman}
\usepackage{amsmath,amssymb,mathtools}
\usepackage{unicode-math}\setmathfont{Latin Modern Math}
\usepackage{xurl,hyperref}
\hypersetup{colorlinks=true,urlcolor=blue}
\setcounter{secnumdepth}{0}\setlength{\parindent}{0pt}\setlength{\parskip}{5pt}\setlength{\emergencystretch}{3em}
\newcommand{\sref}[2]{\hyperlink{source:#1:#2}{[S#2]}}
\newcommand{\cataloguescope}[1]{\paragraph{Recorded target.}#1}
\begin{document}
\section{Kourovka Notebook Problem 21.29}\label{statement}

\subsection{Definitions and mathematical statement}

\paragraph{Shared notation.}$\mathbb N=\{1,2,\ldots\}$, and $\mathbb Z,\mathbb Q,\mathbb R,\mathbb C$ denote the integers, rationals, reals and complex numbers. Any other conventions are specified locally.


Let $G\leq\operatorname{Sym}(\Omega)$ act on a finite set. Write $G_\alpha=\{g\in G:g\alpha=\alpha\}$ and $G_{\alpha,\beta}=G_\alpha\cap G_\beta$ for point and pointwise two-point stabilizers. The action is primitive if it is transitive and has no invariant partition except the singleton and one-part partitions. A regular suborbit is an orbit of a point stabilizer $G_\alpha$ on which that stabilizer acts freely; equivalently, some $G_{\alpha,\beta}$ is trivial.
\paragraph{Statement recorded in the supplied source.}
If $G$ is finite and primitive with a regular suborbit, is it true that
\[
 \forall\alpha,\beta\in\Omega\ \exists\gamma\in\Omega:\quad
 G_{\alpha,\gamma}=G_{\beta,\gamma}=\{1\}?
\]
The source quantifies over all $\alpha,\beta$, including the case $\alpha=\beta$. \sref{statement}{2}
\subsection{Short English statement}
In a finite primitive action that has a trivial two-point stabilizer, can one find a point simultaneously giving trivial stabilizers with any prescribed two points?
\subsection{Sources}
\begin{description}
\item[\hypertarget{source:statement:1}{S1}] ULAM.AI and the UnsolvedMath Contributors, \emph{UnsolvedMath: A Curated Collection of Open Mathematics Problems}, catalogue entry. CC BY 4.0. \url{https://huggingface.co/datasets/ulamai/UnsolvedMath}.
\item[\hypertarget{source:statement:2}{S2}] T. C. Burness and M. Giudici, Problem 21.29, in E. I. Khukhro and V. D. Mazurov (eds.), \emph{The Kourovka Notebook: Unsolved Problems in Group Theory}, 21st edition, arXiv:1401.0300v48 (6 October 2026). This version records a negative answer; the displayed question is the historically registered target. \url{https://arxiv.org/abs/1401.0300v48}.
\end{description}
\label{end:statement}
\end{document}
